\section{The BCS Transformation}\label{sec:bcs}

The BCS transformation~\cite{BCS16} turns an Interactive Oracle Reduction into an Interactive
Reduction by replacing every oracle with a commitment and every oracle query with an opening
argument. It captures both the original vector-commitment/Merkle-tree construction and the
Polynomial IOP + Polynomial Commitment pattern used by Plonk and Marlin.

This section is a blueprint, not a description of existing code:
\texttt{ArkLib/OracleReduction/BCS/Basic.lean} currently contains
\texttt{ProtocolSpec.renameMessage} and two commented-out stubs.

\subsection{The transformation}\label{sec:bcs_transform}

\begin{definition}[Commitment-carrying protocol specification]
    \label{def:bcs_pspec}
    Let $\pSpec : \ProtocolSpec\;n$ with an $\OracleInterface$ on each message. Given a family of
    commitment schemes --- one per message index --- write $\mathsf{Com}_i$ for the commitment type
    of message $i$ and $\pSpec^{\mathsf{open}}_i : \ProtocolSpec\;(n^{\mathsf{open}}_i)$ for the
    opening argument attached to a query on message $i$.

    The transformed specification replaces each message type by its commitment type, using
    \texttt{ProtocolSpec.renameMessage}, and then appends the opening arguments.
    \lean{ProtocolSpec.renameMessage}
\end{definition}

\begin{definition}[BCS transform]
    \label{def:bcs}
    $\mathsf{BCSTransform}$ sends an oracle reduction $R$ to the reduction that
    \begin{enumerate}
        \item runs $R$'s prover, sending a commitment in place of each oracle message;
        \item runs $R$'s oracle verifier to obtain its query list;
        \item runs an opening argument for each query;
        \item accepts iff $R$'s verifier accepts and every opening argument accepts.
    \end{enumerate}
\end{definition}

\subsection{The degree of freedom that decides the formalization}\label{sec:bcs_batching}

The transform's free choices --- the order of the opening arguments, and whether they are batched
--- are usually presented as optimizations. In this library one of them is not: it decides whether
the security proof is reachable at all.

\begin{remark}[Sequential openings inherit the composition blocker]
    \label{rem:bcs_sequential_blocked}
    The natural shape, and the one the commented-out stub sketches, composes the opening arguments
    with $\ProtocolSpec.\mathsf{append}$, giving $n + \sum_i n^{\mathsf{open}}_i$ rounds. Its
    security proof then routes through $\mathsf{Verifier.append\_soundness}$ and its knowledge and
    round-by-round variants.

    Those are unproved, and not for want of effort. $(R_1.\mathsf{append}\;R_2).\mathsf{run}$ runs
    both provers and then both verifiers, whereas running $R_1$ then $R_2$ interleaves
    $P_1, V_1, P_2, V_2$; justifying the swap of $V_1$ past $P_2$ needs the oracle computation
    interpreted into a \emph{commutative} monad, which the VCVio refactor has not yet supplied.

    So a BCS transform built on sequential openings is blocked on a dependency the transform itself
    does not need.
\end{remark}

\begin{remark}[Homomorphic batching avoids it]
    \label{rem:bcs_batched}
    If the commitment scheme is homomorphic, the verifier's queries can be reduced to a single
    opening by taking a random linear combination of the committed messages. The transformed
    protocol is then $\pSpec$ with renamed messages, plus \emph{one} challenge round and
    \emph{one} opening argument --- a fixed, bounded suffix rather than a sum over queries.

    This is worth stating as the primary route rather than the optimization. It removes the
    dependency on sequential-composition security described in
    Remark~\ref{rem:bcs_sequential_blocked}, and it is available for the schemes this library
    already has: Ajtai's commitment is linear, so the combination is immediate.

    The cost is a soundness term for the batching challenge, which must be carried explicitly.
\end{remark}

\subsection{What must be proved}\label{sec:bcs_security}

\begin{theorem}[Completeness]
    \label{thm:bcs_completeness}
    If $R$ is complete with error $\varepsilon$ and every commitment scheme is perfectly correct and
    every opening argument complete, then $\mathsf{BCSTransform}(R)$ is complete with error
    $\varepsilon$.
\end{theorem}

\begin{theorem}[Soundness]
    \label{thm:bcs_soundness}
    If $R$ is sound with error $\varepsilon$, each commitment scheme binding with advantage
    $\delta_i$, and each opening argument sound with error $\eta_i$, then
    $\mathsf{BCSTransform}(R)$ is sound with error $\varepsilon + \sum_i \delta_i + \sum_i \eta_i$,
    plus the batching term of Remark~\ref{rem:bcs_batched} when that route is taken.
\end{theorem}

\begin{theorem}[Knowledge soundness]
    \label{thm:bcs_knowledge_soundness}
    The extractor runs $R$'s extractor against the messages recovered from the opening arguments;
    binding is what makes those messages well defined. Round-by-round variants follow the same
    shape, with the state function extended to record that every commitment opened so far is
    consistent.
\end{theorem}

\subsection{Dependency graph}\label{sec:bcs_deps}

\[
\begin{array}{rcl}
\text{Def.\ref{def:bcs_pspec}} &\to& \text{Def.\ref{def:bcs}} \\
\text{Def.\ref{def:bcs}} + \text{commitment correctness} &\to& \text{Thm.\ref{thm:bcs_completeness}} \\
\text{Def.\ref{def:bcs}} + \text{binding} + \text{opening soundness} &\to& \text{Thm.\ref{thm:bcs_soundness}} \\
\text{Thm.\ref{thm:bcs_soundness}} + \text{$R$'s extractor} &\to& \text{Thm.\ref{thm:bcs_knowledge_soundness}}
\end{array}
\]

Under Remark~\ref{rem:bcs_batched} none of these depends on
$\mathsf{Verifier.append\_soundness}$. Under the sequential route, all three do.

\subsection{Open questions for maintainers}\label{sec:bcs_questions}

\begin{enumerate}
    \item Should the batched form be the definition, with the sequential form derived, or the
        reverse? The argument for batching-first is Remark~\ref{rem:bcs_sequential_blocked}: it is
        the form that can be proved today.
    \item Should $\mathsf{BCSTransform}$ require a homomorphic commitment in its signature, or take
        the batching argument as a parameter so non-homomorphic schemes remain expressible?
    \item Is Merkle capping in scope, or a later specialization?
\end{enumerate}
